\documentclass[aspectratio=169,11pt]{beamer}
\makeatletter
\input{../../templates/dune-professional/beamerthemeDUNEProfessional.sty}
\makeatother
\title{VD detector protection}
\subtitle{System interfaces, proposed interlocks and hardware additions}
\author{Manuel Arroyave}
\date{I\&I workshop draft --- 27 October 2026}
\newcommand{\mydate}{27 October 2026}
\hypersetup{colorlinks=true,urlcolor=DuneInfo,linkcolor=DuneInfo}
\newcommand{\planningnote}[1]{\par\medskip{\footnotesize\color{DuneTextMuted}#1\par}}
% Schematic service boundary, not a surveyed section or controller inventory.
\newcommand{\siteboundarymap}[2]{%
\begin{tikzpicture}[
 box/.style={text=DuneOnAccent,align=center,font=\small,inner sep=4pt},
 level/.style={anchor=east,font=\small,text=DuneTextMuted}]
 \node[font=\small\bfseries,align=center,text width=2.0cm] at (1.25,2.25) {SDSTA\\site services};
 \node[font=\small\bfseries,align=center,text width=2.8cm] at (4.05,2.25) {FFDG DUNE FD\\building services};
 \draw[draw=DuneBorder,line width=1pt] (0.05,1.83) -- (5.6,1.83);
 \draw[draw=DuneSky,line width=2pt] (1.25,1.55) -- (1.25,-1.85);
 \node[level] at (-0.05,1.25) {Surface};
 \node[box,fill=#1,text width=1.95cm] at (1.25,1.25) {Oro Hondo\\Ross shaft};
 \node[level] at (-0.05,0.25) {4850};
 \draw[draw=DuneBorder] (0.05,0.25) -- (5.6,0.25);
 \node[box,fill=#2,text width=2.8cm] at (4.05,0.25) {CUC / caverns\\HVAC, chilled water};
 \node[level] at (-0.05,-0.85) {4910};
 \draw[draw=DuneBorder] (0.05,-0.85) -- (5.6,-0.85);
 \node[box,fill=#2,text width=2.8cm] at (4.05,-0.85) {SP-1 / SP-2\\sump interfaces};
 \node[box,fill=#1,text width=1.95cm] at (1.25,-1.85) {Mine drainage\\and pumping};
 \draw[-{Stealth},draw=DuneSky,line width=1pt] (3.1,-1.3) -- (2.35,-1.65);
\end{tikzpicture}%
}

\begin{document}
\begin{frame}{Definitions and responsibilities}
 \renewcommand{\arraystretch}{1.15}
 \begin{tabularx}{\textwidth}{@{}>{\RaggedRight\arraybackslash}p{2.1cm}>{\RaggedRight\arraybackslash}X@{}}
 \textbf{DPS} & Detector Protection System: arbitrates protection, applies permits and inhibits, and verifies the equipment response.\\
 \textbf{SC} & Slow controls, using Ignition: operates and monitors equipment, coordinates requests, and owns the logs.\\
 \textbf{DAQ} & Configures readout, takes data, and uses detector availability and activity from SC.\\
 \textbf{BMS / cryo} & Building Management Systems and cryogenic controls own their services and provide agreed states to SC.\\
 \textbf{Subsystems} & Define operating limits, required protection and allowed response times.\\
 \textbf{Local\newline interlocks} & Equipment-level protection: can trip without SC or the central DPS.\\
 \end{tabularx}\par\medskip
 \textbf{An operating request cannot override a protection trip.}
 \planningnote{Personnel-safety alarms remain with SDSTA. Agree equipment-protection responsibility for each phase when DPS is not active.}
\end{frame}

\begin{frame}{Facility information reaches DPS through SC}
 \centering
 \begin{tikzpicture}[
 box/.style={fill=DunePaper,text=DuneOnAccent,align=center,text width=3.1cm,minimum height=0.9cm,inner sep=5pt,font=\small},
 edge/.style={-{Stealth},draw=DuneSky,line width=1.2pt},
 duplex/.style={arrows={Stealth-Stealth},draw=DuneSky,line width=1.2pt}]
 \node[box] (ffdg) at (-4.8,2.05) {FFDG DUNE FD BMS\\API via Batavia};
 \node[box] (sdsta) at (0,2.05) {SDSTA BMS\\interface under study};
 \node[box] (cryo) at (4.8,2.05) {Cryogenics in the CUC\\one system to DUNE};
 \node[box,fill=DuneSky,text width=3.2cm] (sc) at (-2.8,0) {SC / Ignition\\operations and logs};
 \node[box,fill=DuneCoral,text width=3.2cm] (dps) at (3.2,0) {DPS PLC\\protection arbitration};
 \draw[edge] (ffdg.south) -- ([xshift=-1.0cm]sc.north);
 \draw[edge,dashed] (sdsta.south) -- (sc.north);
 \draw[edge] (cryo.south) -- ([xshift=1.0cm]sc.north);
 \draw[duplex] (sc.east) -- node[above,font=\small] {OPC UA} (dps.west);
 \draw[draw=DuneBorder,fill=DuneSurface,line width=1pt] (-6.5,-2.55) rectangle (6.5,-1.05);
 \node[font=\small\bfseries] at (0,-1.3) {DETECTOR};
 \node[box,text width=3.4cm,minimum height=0.65cm] (mon) at (-2.8,-1.95) {Subsystems\\control / status};
 \node[box,text width=3.4cm,minimum height=0.65cm] (power) at (3.2,-1.95) {PDUs, supplies\\remote I/O};
 \draw[duplex] (sc.south) -- node[left,font=\small] {OPC UA} (mon.north);
 \draw[duplex] (dps.south) -- node[right,font=\small] {SNMP / EtherCAT} (power.north);
 \end{tikzpicture}\par
 \DuneDiagramCaption{SC connects facility and detector information; DPS decides and applies the protection action.}
 \planningnote{EtherCAT connects remote I/O; wired inhibits connect that I/O to the equipment. SDSTA exchange is under investigation. Proposed central DPS PLC location: CUC.}
\end{frame}

\begin{frame}{FFDG DUNE FD BMS: the Metasys API route}
 \begin{columns}[T,onlytextwidth]
 \begin{column}{0.51\textwidth}
  \centering\siteboundarymap{DunePaper}{DuneSky}
 \end{column}
 \begin{column}{0.47\textwidth}
  \textbf{ISD proposal}\par
  Metasys API in Batavia interfaces the FFDG BMS to DUNE FD SC.\par\medskip
  \textbf{SBN exercise}\par
  The API has been exercised with SBN building data and an OPC UA server for experiments in that building.\par\medskip
  \textbf{Still missing}\par
  The licence for live data. Historical data and OPC UA publication have been exercised.
 \end{column}
 \end{columns}\par\smallskip
 \DuneDiagramCaption{DUNE building services at 4850/4910 connect to wider SURF services. Schematic, not to scale.}
 \planningnote{U-FD-M-901 rev12 and BSI section 230923; API status: Manuel and October test records. Cryogenic process controls are separate from BMS.}
\end{frame}

\begin{frame}{SDSTA BMS: survey the hardware and agree the exchange}
 \begin{columns}[T,onlytextwidth]
 \begin{column}{0.51\textwidth}
  \centering\siteboundarymap{DuneCoral}{DunePaper}
 \end{column}
 \begin{column}{0.47\textwidth}
  \textbf{Hardware survey ongoing}\par
  Identify controllers, I/O, available contacts and network/fibre ports.\par\medskip
  \textbf{Protocols under investigation}\par
  Identify exported states and the path to SC.\par\medskip
  \textbf{Agreement options}\par
  Agree ownership, allowed exchanges, fault reporting and maintenance.
 \end{column}
 \end{columns}\par\smallskip
 \DuneDiagramCaption{SDSTA site services extend beyond the DUNE caverns; BMS connection points remain under survey.}
 \planningnote{4850/4910 maps are ready; access work involves Rachel and SDSTA integrator Dale Curran. Controller models and BMS points remain under survey.}
\end{frame}

\begin{frame}{Cryogenics: one system to SC and DPS}
 \small
 \textbf{Cryogenics is in the CUC.} Its controls are divided internally;
 DUNE receives one coherent system interface.\par\medskip
 \renewcommand{\arraystretch}{1.2}
 \begin{tabularx}{\textwidth}{@{}>{\RaggedRight\arraybackslash}p{3cm}>{\RaggedRight\arraybackslash}X@{}}
 \toprule
 \textbf{Topic} & \textbf{Status and detector consequence}\\\midrule
 Process control & Cryogenics owns regulation and its process/safety controllers; the architecture is documented.\\
 Critical conditions & Low LAr level, relief-valve opening and defined boiling-risk conditions affect HV or CRP-bias permits.\\
 Interface & Cryogenics $\rightarrow$ SC $\leftrightarrow$ DPS. Name the source of each measurement; some level signals belong to detector instrumentation.\\
 Running agreement & Agree readiness states, limits, signal quality, loss response and authorization for each phase.\\
 \bottomrule
 \end{tabularx}
 \planningnote{Cryogenic architecture and 2022 FDR describe the control system. They do not establish an agreed detector run-permit contract. HV/CRP actions come from the interface requirements under review.}
\end{frame}

\begin{frame}{Start with the interlocks proposed by each consortium}
 The matrix collects the current ICD requirements and consortium proposals.
 We review each system against those proposed interlocks.\par\bigskip
 \centering
 \begin{tikzpicture}[
 box/.style={fill=DunePaper,text=DuneOnAccent,align=center,text width=2.2cm,minimum height=1.25cm,inner sep=5pt,font=\small},
 edge/.style={-{Stealth},draw=DuneSky,line width=1.2pt}]
 \node[box] (s) at (-5.5,0) {Condition\\and source};
 \node[box] (i) at (-2.75,0) {Physical input\\and connection};
 \node[box,fill=DuneSky] (c) at (0,0) {DPS / local\\decision};
 \node[box,fill=DuneCoral] (o) at (2.75,0) {Actuator\\and load};
 \node[box] (f) at (5.5,0) {Equipment state\\and feedback};
 \draw[edge] (s)--(i);\draw[edge] (i)--(c);\draw[edge] (c)--(o);\draw[edge] (o)--(f);
 \end{tikzpicture}\par\medskip
 \DuneDiagramCaption{Trace the physical and control path, including the time from detection to completed shutdown.}
 \bigskip\raggedright
 \textbf{For each system:} identify the hardware already provided, missing connections or components, and requirements that change the hardware choice.
 \planningnote{Use the existing rack lists, drawings, parts and power calculations. Count shared I/O once. Proposed actions and unresolved conflicts remain visible in the matrix.}
\end{frame}

\begin{frame}{The DPS brain turns conditions into equipment permits}
 \centering
 \begin{tikzpicture}[
 box/.style={fill=DunePaper,text=DuneOnAccent,align=center,text width=2.65cm,minimum height=1.4cm,inner sep=5pt,font=\small},
 edge/.style={-{Stealth},draw=DuneSky,line width=1.2pt}]
 \node[box] (i) at (-5.1,0) {Conditions, quality\\and SC requests};
 \node[box,fill=DuneSky] (l) at (-1.7,0) {DPS PLC\\rules and trip latches};
 \node[box,fill=DuneCoral] (a) at (1.7,0) {Permits / inhibits\\SNMP /\\remote I/O};
 \node[box] (f) at (5.1,0) {Equipment response\\and feedback};
 \draw[edge] (i)--(l);\draw[edge] (l)--(a);\draw[edge] (a)--(f);
 \end{tikzpicture}\par\medskip
 \DuneDiagramCaption{SC requests an action; the PLC applies the protection rules and reports the result.}
 \bigskip\raggedright
 \textbf{Enable:} request present, required conditions valid, no blocking trip.\par\medskip
 \textbf{Trip:} remove the affected permit, retain the cause and verify the equipment state. SC records the event.\par\medskip
 \textbf{Recovery:} reset and re-enable are separate actions.
 \planningnote{Proposed logic on a commercial PLC based on BDE. Define stale-input and communication-loss responses per load. Local equipment interlocks remain active.}
\end{frame}

\begin{frame}{Detector}
\end{frame}

\begin{frame}{BDE: the existing VD design sets the hardware baseline}
 \centering
 \begin{tikzpicture}[
 box/.style={fill=DunePaper,text=DuneOnAccent,align=center,text width=3.1cm,minimum height=1.05cm,inner sep=6pt},
 edge/.style={-{Stealth},draw=DuneSky,line width=1.2pt}]
 \node[box,fill=DuneSky] (io) at (-4,0) {BDE auxiliary rack\\remote I/O};
 \node[box] (lv) at (1.6,1.1) {2 LV racks\\2 PL506 supplies};
 \node[box,fill=DuneCoral] (wiec) at (1.6,-1.1) {8 WIECs\\8 PTCs + 48 WIBs};
 \draw[edge] (io.north east) -| (-1.7,1.1) -- node[above,font=\small] {inhibits} (lv.west);
 \draw[arrows={Stealth-Stealth},draw=DuneSky,line width=1.2pt] (io.south east) -| (-1.7,-1.1) -- node[below,font=\small] {EtherCAT} (wiec.west);
 \draw[edge] (lv.east) -- (4.35,1.1) -- node[right,font=\small] {48 V} (4.35,-1.1) -- (wiec.east);
 \end{tikzpicture}\par\medskip
 \DuneDiagramCaption{Ten groups: 10 auxiliary racks, 20 LV racks, 20 PL506 supplies, 80 WIECs and 480 WIBs.}
 \medskip\raggedright
 The panel drawings, parts lists, power calculations and Ignition screens are available.
 Ten bias racks provide the separate CRP/FC protection interfaces.
 \planningnote{VD rack workbook and June SC layout. BDE = bottom drift electronics; WIEC = warm interface electronics crate; PTC = power and timing card; WIB = warm interface board. Design/test evidence is not a claim of VD installation.}
\end{frame}

\begin{frame}{BDE trip example: fan cooling and the bias dependency}
 \centering
 \begin{tikzpicture}[
 box/.style={fill=DunePaper,text=DuneOnAccent,align=center,text width=2.35cm,minimum height=1.05cm,inner sep=5pt,font=\small},
 edge/.style={-{Stealth},draw=DuneSky,line width=1.2pt}]
 \node[box] (fan) at (-5.1,0) {Fan status\\DB25 / EL1004};
 \node[box,fill=DuneSky] (plc) at (-1.7,0) {DPS PLC\\cooling rule};
 \node[box,fill=DuneCoral] (lv) at (1.7,1) {EL2124 output\\PL506 inhibit};
 \node[box] (wiec) at (5.1,1) {Affected WIEC\\48 V off};
 \node[box,fill=DuneCoral] (bias) at (1.7,-1) {Bias remote I/O\\ISEG inhibit};
 \node[box] (crp) at (5.1,-1) {Associated\\CRP bias off};
 \draw[edge] (fan)--node[above=18pt,font=\small] {EtherCAT}(plc);
 \draw[edge] (plc.north east) |- (lv.west);\draw[edge] (lv)--(wiec);
 \draw[edge,dashed] (plc.south east) |- (bias.west);\draw[edge,dashed] (bias)--(crp);
 \end{tikzpicture}\par\medskip
 \DuneDiagramCaption{Solid: documented BDE cooling path. Dashed: associated bias inhibition, if required by the VD rule.}
 \medskip\raggedright
 The output paths already exist. Map the affected WIEC and bias channels, set the cooling rule, and verify supply voltage/state after the trip.
 \planningnote{The historical HD rule links cooling and bias. The VD cooling proposal currently assigns LV inhibition; its bias dependency needs agreement. Six fan-status inputs per WIEC imply 48 per eight-WIEC group, not 48 new inputs.}
\end{frame}

\begin{frame}{PDS DAPHNE: control, power and the DAQ interface}
 \centering
 \begin{tikzpicture}[
 box/.style={fill=DunePaper,text=DuneOnAccent,align=center,text width=2.4cm,minimum height=0.8cm,inner sep=5pt,font=\small},
 edge/.style={-{Stealth},draw=DuneSky,line width=1.2pt},
 duplex/.style={arrows={Stealth-Stealth},draw=DuneSky,line width=1.2pt}]
 \node[box,fill=DuneSky] (dps) at (-5,1.15) {DPS PLC};
 \node[box] (sc) at (-1.2,1.15) {SC / Ignition};
 \node[box] (daq) at (3.8,1.15) {DAQ\\CCM / Hermes};
 \node[box,text width=3.0cm] (ua) at (-1.2,-0.2) {PDS OPC UA interface};
 \node[box,fill=DuneCoral] (lv) at (-5,-1.45) {WIENER LV PS};
 \node[box,text width=3.0cm] (daphne) at (3.8,-1.45) {DAPHNE\\bias control / readout};
 \draw[duplex] (dps)--node[above,font=\small]{OPC UA}(sc);
 \draw[duplex] (sc)--node[left,font=\small]{OPC UA}(ua);
 \draw[duplex] (daq.south west) |- node[above,pos=0.75,font=\small]{OPC UA}(ua.east);
 \draw[duplex] (ua.south) |- node[above,pos=0.7,font=\small]{ZMQ control / status}(daphne.west);
 \draw[duplex] (daq.south) -- node[right,font=\small,align=left]{IPBus / UDP\\configuration}(daphne.north);
 \draw[duplex] (dps)--node[left,font=\small]{SNMP}(lv);
 \draw[edge,draw=DuneCoral] (lv.south)--(-5,-2.15)--node[below,font=\small]{48 V power}(3.8,-2.15)--(daphne.south);
 \end{tikzpicture}\par\medskip
 \DuneDiagramCaption{Working ICD DAPHNE paths plus DPS and LV control. PDS power scope: 34 LV feeds.}
 \medskip\raggedright
 \textbf{SiPM-bias inhibit:} DPS $\rightarrow$ SC $\rightarrow$ DAPHNE. Bias state returns to SC; DAQ receives detector availability.
 \planningnote{DAPHNE-controlled bias only; cathode power follows PoF. Test inhibit acknowledgement and timing.}
\end{frame}

\begin{frame}{PDS PoF: monitor locally, inhibit through the PDU}
 \centering
 \begin{tikzpicture}[
 box/.style={fill=DunePaper,text=DuneOnAccent,align=center,text width=2.7cm,minimum height=0.85cm,inner sep=5pt,font=\small},
 edge/.style={-{Stealth},draw=DuneSky,line width=1.2pt},
 duplex/.style={arrows={Stealth-Stealth},draw=DuneSky,line width=1.2pt}]
 \node[box] (sc) at (-4.8,1) {SC / Ignition};
 \node[box] (local) at (0,1) {Opta PLC /\\Beckhoff PLC\\local monitoring};
 \node[box,fill=DuneCoral] (box) at (4.8,0) {PoF laser box\\4 voltage readings\\1 temperature};
 \node[box,fill=DuneSky] (dps) at (-4.8,-1.2) {DPS PLC};
 \node[box] (pdu) at (0,-1.2) {Rack PDU};
 \draw[duplex] (sc)--node[above,font=\small]{OPC UA}(local);
 \draw[edge] (box.north west) |- node[above,pos=0.75,font=\small]{analog inputs}(local.east);
 \draw[duplex] (sc)--node[left,font=\small]{OPC UA}(dps);
 \draw[duplex] (dps)--node[above,font=\small]{SNMP}(pdu);
 \draw[edge,draw=DuneCoral] (pdu.east) -| node[below,pos=0.25,font=\small]{AC power}(box.south);
 \draw[edge,dashed] (local.south east)--node[below,sloped,font=\small]{optional enable}(box.west);
 \end{tikzpicture}\par\medskip
 \DuneDiagramCaption{Selected inhibition route: switch the mapped PDU outlet. Explore a separate box-enable input if available.}
 \medskip\raggedright
 \textbf{80 boxes: 320 voltage + 80 temperature readings = 400 readings.}
 Size each controller option for these physical inputs and its box grouping.
 \planningnote{Four voltage readings and one chassis temperature per box: Manuel's current scope. Readings do not count laser drivers or central DPS terminals. The optional control signal and feedback remain to be identified.}
\end{frame}

\begin{frame}{HV: drift voltage and CRP / field-cage bias}
 \textbf{VD inventory:} one Heinzinger drift-HV supply; ten ISEG bias crates,
 576 channels = 480 CRP + 96 field-cage channels.\par\medskip
 \renewcommand{\arraystretch}{1.35}
 \begin{tabularx}{\textwidth}{@{}>{\RaggedRight\arraybackslash}p{2.2cm}>{\RaggedRight\arraybackslash}X>{\RaggedRight\arraybackslash}X@{}}
 \toprule
 \textbf{Equipment} & \textbf{Control / monitoring} & \textbf{Protection path}\\\midrule
 Drift HV & SC / OPC UA $\leftrightarrow$ LabVIEW / cRIO; analog setpoint and V/I readback & DPS $\rightarrow$ EtherCAT remote I/O $\rightarrow$ wired Heinzinger interlock\\
 CRP / FC bias & SNMP supply control/status; SC receives states and logs through DPS & DPS $\rightarrow$ EtherCAT remote I/O $\rightarrow$ wired ISEG channel inhibits\\
 \bottomrule
 \end{tabularx}\par\medskip
 \textbf{Example:} low LAr level removes the drift-HV permit. A drift-HV trip also shuts down the associated field-cage termination supplies.
 \planningnote{EDMS 3315615 v1 attachment, internal v2; October architecture. Map actual pins, channels and feedback before counting added I/O. Supply shutoff and voltage decay are separate timing requirements.}
\end{frame}

\begin{frame}{CalCI temperature sensors and their readout}
 \centering
 \begin{tikzpicture}[
 box/.style={fill=DunePaper,text=DuneOnAccent,align=center,text width=2.5cm,minimum height=0.85cm,inner sep=5pt,font=\small},
 edge/.style={-{Stealth},draw=DuneSky,line width=1.2pt}]
 \node[box] (rtd) at (-5,1) {RTD sensors\\warm cables};
 \node[box] (crio) at (-1.5,1) {Readout / cRIO\\LabVIEW};
 \node[box] (fbg) at (-5,-1) {Fibre temperature\\sensors};
 \node[box] (int) at (-1.5,-1) {Interrogator\\readout service};
 \node[box] (sc) at (2.5,0) {SC / OPC UA\\logs and quality};
 \node[box,fill=DuneSky,text width=1.5cm] (dps) at (5.8,0) {DPS};
 \draw[edge] (rtd)--(crio);\draw[edge] (fbg)--(int);
 \draw[edge] (crio.east)-| (sc.north);\draw[edge] (int.east)-| (sc.south);
 \draw[edge] (sc)--(dps);
 \end{tikzpicture}\par\medskip
 \DuneDiagramCaption{Detector instrumentation is separate from cryogenic process control. Name the sensor and its physical location.}
 \medskip\raggedright
 Temperature profiles support cooldown and operation. A protection rule needs the affected equipment, threshold, persistence and allowed response time.\par\medskip
 \textbf{Hardware review:} match RTD / fibre channels, warm readout and service interfaces to the required measurements.
 \planningnote{October SC architecture and CalCI inventory. Trace level measurements to their own instrument and readout.}
\end{frame}

\begin{frame}{CalCI lasers: coordinate emission and detector operation}
 \centering
 \begin{tikzpicture}[
 box/.style={fill=DunePaper,text=DuneOnAccent,align=center,text width=2.6cm,minimum height=0.9cm,inner sep=5pt,font=\small},
 edge/.style={-{Stealth},draw=DuneSky,line width=1.2pt},
 duplex/.style={arrows={Stealth-Stealth},draw=DuneSky,line width=1.2pt}]
 \node[box] (sc) at (-4.9,0.7) {SC / OPC UA};
 \node[box] (cib) at (0,0.7) {Calibration interface\\board / controller};
 \node[box,fill=DuneCoral] (laser) at (4.9,0.7) {Ionization laser\\station};
 \node[box] (daq) at (-4.9,-1.2) {DAQ / calibration\\request and activity};
 \node[box,fill=DuneSky] (dps) at (0,-1.2) {DPS PLC\\via SC};
 \node[box] (local) at (4.9,-1.2) {Local laser\\safety interlocks};
 \draw[duplex] (sc)--(cib);\draw[duplex] (cib)--(laser);
 \draw[duplex] (daq)--(sc);\draw[edge] (local)--(laser);
 \draw[duplex] (dps.north west)--node[above,sloped,font=\small]{OPC UA}(sc.south east);
 \end{tikzpicture}\par\medskip
 \DuneDiagramCaption{Working inventory: four ionization-laser stations. The CIB has an OPC UA interface in the SC architecture.}
 \medskip\raggedright
 \textbf{Proposed operating rule:} deny conflicting laser/PDS requests. Require emission or shutter-off confirmation before PDS return.\par\medskip
 \textbf{Hardware review:} station power, permits, shutter/emission feedback and trigger connection.
 \planningnote{CalCI inventory; October SC architecture; proposed rule O-002. Local personnel-safety interlocks stay independent.}
\end{frame}

\begin{frame}{CalCI purity monitors: separate activity from protection}
 \centering
 \begin{tikzpicture}[
 box/.style={fill=DunePaper,text=DuneOnAccent,align=center,text width=2.7cm,minimum height=1.0cm,inner sep=5pt,font=\small},
 edge/.style={-{Stealth},draw=DuneSky,line width=1.2pt},
 duplex/.style={arrows={Stealth-Stealth},draw=DuneSky,line width=1.2pt}]
 \node[box] (prm) at (-5.1,0) {PrM bias / front-end\\xenon flash lamp};
 \node[box] (read) at (-1.7,0) {Digitizer / LabVIEW\\control and readout};
 \node[box] (sc) at (1.7,0) {SC / OPC UA\\activity and status};
 \node[box] (daq) at (5.1,0) {DAQ\\activity tags};
 \draw[duplex] (prm)--(read);\draw[duplex] (read)--(sc);\draw[edge] (sc)--(daq);
 \end{tikzpicture}\par\medskip
 \DuneDiagramCaption{Track lamp activity, PrM supply state and the affected detector region.}
 \bigskip\raggedright
 \textbf{Proposal Q-001:} allow acquisition and publish the activity interval for DAQ tagging.\par\medskip
 \textbf{Open ICD conflict:} the older PDS ICD requests SiPM bias off while the purity monitor is active. The consortium must agree which rule applies.\par\medskip
 \textbf{Hardware review:} lamp control/status, HV/LV supplies, digitizer interfaces and the readback used to verify the action.
 \planningnote{Working PrM inventory; EDMS 3309688 v1 section 5; SC--PDS--DPS working matrix. No automatic bias trip is assigned by Q-001.}
\end{frame}

\begin{frame}{Racks: one protection package shared by the subsystems}
  \centering
  \begin{tikzpicture}[
    box/.style={fill=DunePaper,text=DuneOnAccent,align=center,
      text width=2.0cm,minimum height=1.05cm,inner sep=5pt,font=\small},
    edge/.style={-{Stealth},draw=DuneSky,line width=1.2pt},
    duplex/.style={arrows={Stealth-Stealth},draw=DuneSky,line width=1.2pt}]
    \draw[draw=DuneBorder,line width=1pt,fill=DuneSurface] (0.3,-0.9) rectangle (3.7,2.45);
    \node[font=\small\bfseries] at (2,2.2) {Rack infrastructure};
    \node[box,text width=2.2cm] (sc) at (-5.5,0) {SC / Ignition\\operations\\and logs};
    \node[box,fill=DuneSky] (plc) at (-1.8,0) {DPS PLC\\arbitration};
    \node[box,fill=DuneCoral,text width=2.4cm] (pdu) at (2,0) {PDU\\local trips\\outlet control};
    \node[box,text width=1.8cm] (load) at (5.6,0) {Subsystem\\loads};
    \node[box,text width=2.8cm,minimum height=0.9cm] (sensors) at (2,1.45) {Smoke + humidity\\temperature / leak};
    \draw[duplex] (sc.east) -- node[above,font=\small] {OPC UA} (plc.west);
    \draw[duplex] (plc.east) -- node[above,font=\small] {SNMP} (pdu.west);
    \draw[edge] (pdu.east) -- node[above,font=\small] {AC power} (load.west);
    \draw[edge] (sensors.south) -- (pdu.north);
  \end{tikzpicture}\par
  \medskip
  \DuneDiagramCaption{Count the rack, PDUs, sensors and interface parts once. Link each subsystem to its PDU outlets.}
  \medskip
  \raggedright Local PDU rules protect the mapped outlets. Humidity is monitored;
  a humidity shutdown needs a stated requirement.\par
  \planningnote{Rack inventory includes 20 TDE racks: current DPS control is through their PDUs. Crate status and external inhibit remain future options. Reset and re-enable are separate actions.}
\end{frame}

\begin{frame}{Test arbitration, trips and recovery on the bench}
  \centering
  \begin{tikzpicture}[
    box/.style={fill=DunePaper,text=DuneOnAccent,align=center,
      text width=2.25cm,minimum height=1.0cm,inner sep=5pt},
    edge/.style={-{Stealth},draw=DuneSky,line width=1.2pt}]
    \node[box] (in) at (-4.8,0) {Fault input or\\conflicting request};
    \node[box,fill=DuneSky] (plc) at (-1.6,0) {DPS PLC\\arbitration};
    \node[box,fill=DuneCoral] (out) at (1.6,0) {Supply inhibit\\or PDU command};
    \node[box] (state) at (4.8,0) {Equipment state\\and feedback};
    \draw[edge] (in) -- (plc);
    \draw[edge] (plc) -- (out);
    \draw[edge] (out) -- (state);
  \end{tikzpicture}\par
  \medskip
  \DuneDiagramCaption{Use the BDE vertical-slice work and G271 equipment for representative VD paths.}
  \bigskip
  Measure detection through completed shutdown. Test local trips, SC and fieldbus
  loss, reset and re-enable. Verify that SC logs the decisions and equipment states.\par
  \bigskip
  No critical sub-ms function has been identified so far. Simulation prepares the
  fault and load cases; hardware tests measure the complete response.
\end{frame}

\begin{frame}{Take the VD hardware to review in June 2027}
  \renewcommand{\arraystretch}{1.4}
  \begin{tabularx}{\textwidth}{@{}>{\RaggedRight\arraybackslash}p{3.2cm}>{\RaggedRight\arraybackslash}X@{}}
    \toprule
    \textbf{Work before review} & \textbf{Result for procurement}\\
    \midrule
    Trace the VD inhibits & Existing hardware, added channels and cables, quantities and open alternatives\\
    Build representative paths & Measured trip time, failure/recovery results and SC logging\\
    Prepare purchases & Panel space and power, drawings, parts lists, costs, lead times and spares\\
    \bottomrule
  \end{tabularx}\par
  \bigskip
  After readiness review: order the hardware and install it in the site sequence.
  \planningnote{Possible facility-interface need: summer 2027. Full VD response-matrix review: early 2028. The current path review establishes the requirements that affect hardware purchases.}
\end{frame}

\begin{frame}{I will lead the VD hardware work with a control engineer}
  \renewcommand{\arraystretch}{1.4}
  \begin{tabularx}{\textwidth}{@{}>{\RaggedRight\arraybackslash}p{4.1cm}>{\RaggedRight\arraybackslash}X@{}}
    \toprule
    \textbf{Proposed contribution} & \textbf{Work}\\
    \midrule
    Manuel + control engineer & VD design, SURF interfaces and procurement package\\
    Terri and Trevor & BDE panels, PLC/I/O implementation and bench work\\
    Geoff + Manuel & SC commands, PDU protection and recovery\\
    Vivian + consortia / DPS & Equipment requirements and ICD updates\\
    \bottomrule
  \end{tabularx}\par
  \bigskip
  \textbf{Support needed:} control-engineering time and access to the existing test stands.
  \planningnote{SURF work: inspection maps are ready; access is being coordinated with Rachel, with Dale Curran identified as an SDSTA integration contact.}
\end{frame}

\appendix
\begin{frame}{BDE drawings, parts and software are already available}
  \begin{columns}[T,onlytextwidth]
    \begin{column}{0.47\textwidth}
      \textbf{TPC protection panel designs}\par\medskip
      \colorbox{DunePaper}{\includegraphics[page=2,width=\dimexpr\linewidth-2\fboxsep\relax,height=3.1cm,keepaspectratio,
        trim=290bp 500bp 900bp 430bp,clip]{figures/tpc-panel-reference.pdf}}\par
      \medskip\DuneDiagramCaption{Central PLC, bias I/O and LV/fan/heater reference panels.}
    \end{column}
    \begin{column}{0.49\textwidth}
      \textbf{BDE Ignition screens}\par\medskip
      \includegraphics[page=7,width=\linewidth,height=3.1cm,keepaspectratio,
        trim=10bp 210bp 460bp 74bp,clip]{figures/sc-progress-june2026.pdf}\par
      \planningnote{Existing display of WIEC groups; June 2026 SC progress presentation.}
    \end{column}
  \end{columns}\par
  \medskip
  Parts list: CX5140 PLC, EK1100 couplers, EK1521 fibre links and EL2124 inhibit outputs.
  \planningnote{Historical TPC reference: EDMS 2401090 v4, sheet 2, with parts and power workbooks. The VD baseline is eight WIECs per group.}
\end{frame}

\begin{frame}{BDE cooling already has an input-to-inhibit path}
  \centering
  \begin{tikzpicture}[
    box/.style={fill=DunePaper,text=DuneOnAccent,align=center,
      text width=1.75cm,minimum height=1.25cm,inner sep=4pt,font=\small},
    edge/.style={-{Stealth},draw=DuneSky,line width=1.2pt}]
    \node[box] (fan) at (-5.75,0.6) {Fan status\\DB25 cable};
    \node[box] (input) at (-3.45,0.6) {EL1004\\input I/O};
    \node[box,fill=DuneSky] (plc) at (-1.15,0.6) {DPS PLC\\cooling rule};
    \node[box] (output) at (1.15,0.6) {EL2124\\5 V output};
    \node[box,fill=DuneCoral] (supply) at (3.45,0.6) {PL506\\channel inhibit};
    \node[box] (load) at (5.75,0.6) {Affected WIEC\\48 V feed};
    \draw[edge] (fan) -- (input);
    \draw[edge] (input) -- (plc);
    \draw[edge] (plc) -- (output);
    \draw[edge] (output) -- (supply);
    \draw[edge] (supply) -- (load);
    \node[font=\small] at (-2.3,1.55) {EtherCAT};
    \node[font=\small] at (0,1.55) {EtherCAT};
    \node[font=\small] at (2.3,1.55) {Wired};
  \end{tikzpicture}\par
  \medskip
  \DuneDiagramCaption{Cooling fault $\rightarrow$ withdraw the permit for the affected PL506 channel.}
  \medskip
  \raggedright \textbf{Feedback:} PL506 state and voltage/current $\rightarrow$ SNMP $\rightarrow$ DPS;
  OPC UA carries the result to SC for logging.\par
  \medskip
  \textbf{VD allocation check:} six fan-status signals per WIEC means 48 per
  eight-WIEC group. Match these to the populated VD inputs and cables.\par
  \planningnote{The six-fan WIEC note differs from the four-fan ICEBERG reference. These are baseline requirements, not new-module counts. Fan voting and allowed shutdown time remain to be agreed.}
\end{frame}

\begin{frame}{Backup: DAPHNE paths in the working ICD}
 \centering
 \resizebox{0.84\linewidth}{!}{% Adapted from icd-sc-pds/sections/03_protocols_services.tex at bfb2a8b.
\begin{tikzpicture}[
 box/.style={fill=DunePaper,text=DuneOnAccent,align=center,font=\small,inner sep=5pt},
 flow/.style={-{Stealth},draw=DuneSky,line width=1.2pt},
 both/.style={arrows={Stealth-Stealth},draw=DuneSky,line width=1.2pt}]
\node[box,text width=37mm] (ignition) at (0,0)
 {SC server\\OPC UA client\\$\rightarrow$ Ignition};
\node[box,text width=32mm] (daq) at (5.0,0) {DAQ CCM client\\opmon subscriber};
\node[box,text width=35mm] (hermesclient) at (10.7,0)
 {DAQ Hermes\\configuration client};
\node[box,text width=75mm,minimum height=19mm] (ua) at (2.5,-2.3)
 {SC bridge server\\ZMQ client $\rightarrow$ OPC UA bridge\\Integrated DAQ configuration};
\node[box,text width=49mm,minimum height=18mm] (daphne) at (2.5,-5.2)
 {DAPHNE\\AXI / OS $\rightarrow$ daphne-server\\$\rightarrow$ ZMQ};
\node[box,text width=37mm,minimum height=18mm] (hermes) at (10.7,-5.2)
 {Board-side Hermes\\IPBus / UDP\\Configuration endpoint};
\node[draw=DuneSky,dashed,rounded corners,fit=(daphne)(hermes),inner sep=8pt,
 label={[font=\small,text=DuneSky]below:One DAPHNE board; one shared 1\,Gb Ethernet CCM connection}] {};
\draw[both] (ignition) -- node[left,font=\small]{OPC UA} (ignition |- ua.north);
\draw[both] (daq) -- node[right,font=\small]{OPC UA} (daq |- ua.north);
\draw[both] ([xshift=-11mm]ua.south) -- node[left,align=right,font=\small,fill=DuneBackground,inner sep=2pt]
 {Configuration\\request/response} ([xshift=-11mm]daphne.north);
\draw[flow,draw=DuneSky] ([xshift=11mm]daphne.north) -- node[right,align=left,font=\small,fill=DuneBackground,inner sep=2pt]
 {Published\\monitoring} ([xshift=11mm]ua.south);
\draw[both] (hermesclient.south) -- node[right,align=left,font=\small]
 {IPBus / UDP\\configuration} (hermes.north);
\end{tikzpicture}
}\par
 \DuneDiagramCaption{Working SC--PDS--DPS ICD, Figure 1. DAQ CCM/opmon and Hermes remain distinct clients.}
\end{frame}

\begin{frame}{Backup: VD rack space and the existing DPS allocation}
  \renewcommand{\arraystretch}{1.4}
  \begin{tabularx}{\textwidth}{@{}>{\RaggedRight\arraybackslash}p{3.4cm}r>{\RaggedRight\arraybackslash}X@{}}
    \toprule
    \textbf{Rack type} & \textbf{Count} & \textbf{DPS allocation and workbook space}\\
    \midrule
    BDE auxiliary & 10 & 12U DPS; U3--U21 marked Space (19U); eight WIECs per rack\\
    BDE low voltage & 20 & One PL506 per rack; four WIECs per supply\\
    Bias high voltage & 10 & 12U DPS; U3--U12 marked Space (10U)\\
    \bottomrule
  \end{tabularx}\par
  \bigskip
  Space marked in the rack workbook is not a new DPS reservation.
  Free DIN-rail length and I/O capacity depend on the populated VD panel.
  \planningnote{May 2025 VD workbook: BDE Aux Rack and Bias HV Rack. Auxiliary U9 also has a possible ground-plane readout note. U is rack height.}
\end{frame}

\begin{frame}{The VD bias-panel draft specifies two DIN rails per rack}
  \renewcommand{\arraystretch}{1.3}
  \begin{tabularx}{\textwidth}{@{}>{\RaggedRight\arraybackslash}Xrr@{}}
    \toprule
    \textbf{Component} & \textbf{Per rail} & \textbf{10 racks}\\
    \midrule
    EK1100 EtherCAT coupler & 1 & 20\\
    EL9410 supply terminal & 1 & 20\\
    EL2124: four 5 V outputs & 18 & 360\\
    EL9011 end cap & 1 & 20\\
    24 V power supply & 1 & 20\\
    \bottomrule
  \end{tabularx}\par
  \medskip
  Each rail occupies 337 mm. The draft provides 1,440 output channels;
  the working VD inventory lists 576 bias channels.\par
  \planningnote{May 2025 workbook, Bias HV Rack!K17:K21. This is the first-pass layout; the difference from 576 bias channels is not verified spare capacity.}
\end{frame}

\begin{frame}{Use the component loads for the VD power budget}
  \small
  \renewcommand{\arraystretch}{1.3}
  \begin{tabularx}{\textwidth}{@{}>{\RaggedRight\arraybackslash}p{3.7cm}>{\RaggedRight\arraybackslash}p{2.4cm}>{\RaggedRight\arraybackslash}X@{}}
    \toprule
    \textbf{Equipment} & \textbf{Power / capacity} & \textbf{Basis and limit}\\
    \midrule
    PL506 input & 2.25 kW each & 20 supplies: about 45.1 kW expected AC input\\
    WIEC fans only & 380 W per group & $8\times6\times0.33\,\mathrm{A}\times24\,\mathrm{V}$; excludes heaters and I/O\\
    Bias DPS supply allowance & 24 V, 40 A & 960 W DC capacity per rack in the workbook\\
    \bottomrule
  \end{tabularx}\par
  \medskip
  The inherited 831 W auxiliary figure needs a VD component budget
  covering fans, heaters and populated I/O.\par
  \planningnote{May workbook; WIEC safety note, internal revision 0. Fan arithmetic assumes six fans per WIEC. Rated load, expected AC input and DC supply capacity are distinct.}
\end{frame}

\begin{frame}{Rack infrastructure already has a hardware reference}
  \small
  \renewcommand{\arraystretch}{1.2}
  \begin{tabularx}{\textwidth}{@{}>{\RaggedRight\arraybackslash}p{3.4cm}>{\RaggedRight\arraybackslash}X@{}}
    \toprule
    \textbf{Rack component} & \textbf{Reference hardware / allocation}\\
    \midrule
    PDU & Raritan PX3; inlet and outlet allocation for the connected loads\\
    Smoke detector & 4WTAR-B and EOLR-1 end-of-line relay\\
    Contact interface & DX-PD2C5; 12 V supply to the smoke detector\\
    Humidity & Record the probe and peripheral interface for each rack\\
    Temperature / water leak & PDU peripheral sensors where fitted\\
    \bottomrule
  \end{tabularx}\par
  \medskip
  \normalsize The May workbook includes a 76-entry rack-protection/PDU allowance.
  Map the actual rack packages against this existing allowance.\par
  \planningnote{DUNE-doc-32655 v8, sections 3 and 5; May workbook, Standard Rack and Power Summary!B28. The guide does not specify a humidity probe model or trip requirement.}
\end{frame}

\begin{frame}{Reference: local PDU trips and recovery}
  \renewcommand{\arraystretch}{1.45}
  \begin{tabularx}{\textwidth}{@{}>{\RaggedRight\arraybackslash}p{3.7cm}>{\RaggedRight\arraybackslash}X@{}}
    \toprule
    \textbf{PX3 setup guide} & \textbf{Local behaviour documented}\\
    \midrule
    Smoke, water, temperature & Local event rules and scripts can switch off the protected outlets\\
    Network-switch outlet & Kept powered during the guide's fallback shutdown\\
    After a local trip & Power is restored deliberately; no automatic restoration\\
    \bottomrule
  \end{tabularx}\par
  \bigskip
  The VD control chain is Ignition $\rightarrow$ DPS PLC $\rightarrow$ PDU.
  Local trip status returns to the PLC and to SC for logging.
  \planningnote{DUNE-doc-32655 v8 describes an Ignition-direct prototype with local fallback on loss of Ignition. The VD PLC command path requires matching reachability, trip and reset configuration; it is not a claim that the prototype already implements it.}
\end{frame}

\begin{frame}{The VD inventory covers the other subsystem connections}
  \renewcommand{\arraystretch}{1.35}
  \begin{tabularx}{\textwidth}{@{}>{\RaggedRight\arraybackslash}p{3.2cm}>{\RaggedRight\arraybackslash}X@{}}
    \toprule
    \textbf{Subsystem} & \textbf{Working inventory and interfaces}\\
    \midrule
    PDS readout / calibration & 34 LV feeds in the current presentation scope; supply status and control interfaces\\
    Power over fibre (PoF) & 80 boxes; 320 voltage and 80 temperature readings; Opta PLC or Beckhoff PLC alternative\\
    TDE & 20 mezzanine racks; PDU control now; other TDE interfaces are possible extensions\\
    Calibration / cryogenic instrumentation & Level and temperature readout, purity monitors, ionization laser and neutron source\\
    \bottomrule
  \end{tabularx}
  \planningnote{Subsystem exchanges use OPC UA through SC. PoF uses the current 4 V + 1 temperature case; these readings do not count central DPS terminals.}
\end{frame}

\begin{frame}{The BMS drawing identifies 20 control-panel symbols}
  \renewcommand{\arraystretch}{1.4}
  \begin{tabularx}{\textwidth}{@{}>{\RaggedRight\arraybackslash}p{3.4cm}>{\RaggedRight\arraybackslash}X@{}}
    \toprule
    \textbf{Drawing area} & \textbf{Interfaces on the inspection maps}\\
    \midrule
    CUC, level 4850 & Eight panel references: air handling, chilled water, power and associated equipment\\
    Level 4910 & North and south sump interfaces, SP-1 and SP-2\\
    4850-08-ES group & Spare panel for future equipment; mapped to the access-drift inspection area\\
    Other 4850 areas & Chillers, pumps, ventilation and generator interfaces\\
    \bottomrule
  \end{tabularx}\par
  \medskip
  Inspection planned with Rachel; Dale Curran is the SDSTA integration contact identified for this work.
  \planningnote{U-FD-M-901 rev12, 14 February 2025; marked 4850/4910 maps and survey register. Symbols are drawing references; field observations are pending.}
\end{frame}

\begin{frame}{Reference: the May 2025 VD rack inventory}
  \renewcommand{\arraystretch}{1.15}
  \begin{tabularx}{\textwidth}{@{}>{\RaggedRight\arraybackslash}Xr>{\RaggedRight\arraybackslash}Xr@{}}
    \toprule
    \textbf{Rack type} & \textbf{Count} & \textbf{Rack type} & \textbf{Count}\\
    \midrule
    BDE low voltage & 20 & PDS low voltage & 4\\
    BDE auxiliary & 10 & Drift HV & 2\\
    TDE & 20 & Reserved drift HV & 2\\
    Bias HV & 10 & Spare / CalCI & 8\\
    \bottomrule
  \end{tabularx}\par
  \bigskip
  \textbf{76 numbered mezzanine positions:} SEDM-01 to SEDM-76.\par
  \medskip
  The worksheet also has a separate concept for 40 top-plate mini-racks.
  \planningnote{May 2025 allocation: rack positions, distinct from PLC panels and remote-I/O stations. Includes the listed reserved and spare positions.}
\end{frame}

\end{document}
